\documentclass[12pt]{book}
\usepackage{amsmath}
\usepackage{amssymb}
\usepackage{geometry}
\geometry{margin=1in}

\title{Volume 13: Philosophy, Methodology, and Programmatics \\ Book 02: Investigations and Prompts}
\author{James C. Harvey}
\date{December 2025}

\begin{document}
\maketitle
\tableofcontents

\chapter{Investigations as Formal Extensions}

\section*{Abstract}
Investigations extend SDT into specific empirical or mathematical tests. Each investigation is a
structured application of SDT operators, and the prompts define the scope of those applications.

\section{Investigation Protocol}
Every investigation must state the geometry, define occlusion conditions, and declare validation criteria.

\chapter{Prompt Architecture}

\section*{Abstract}
Prompts define the investigation’s boundary and expected outputs. They are formal constraints, not
informal requests.

\section{Prompt Structure}
Each prompt specifies: geometry inputs, required operators, predicted outputs, and validation sources.

\chapter{Investigation Templates}

\section*{Abstract}
Templates ensure that investigations are reproducible and compatible with benchmark standards.

\section{Template Fields}
Templates include: constants, state vector construction, occlusion evaluation, regime selection, and
result comparison.
\end{document}
